% Folder 74, batch 03 — archivists' pages 41–60.
% Pass: Opus 5.5 (claude-opus-5-5), 2026-10-03 — first pass, unchecked against the pages by a human.
%
% Skipped: p. 42 and p. 45 (blank covers carrying only a pencilled number);
% the upper half of p. 41 (household accounts and an address).
% Recurring doubtful term: « W-liées » (see the note on p. 46).
\documentclass[11pt,a4paper]{article}
\input{../preamble/grothendieck}

\folder{74}
\batch{3}
\pages{41}{60}
\foldertitle{Complexes cubiques : notes manuscrites (s.d.).}
\dating{[à partir de 1976]}
\watermark{Édition de démonstration}

\begin{document}

\page{41}
\note{la moitié inférieure du feuillet, écrite tête-bêche ; la moitié supérieure ne porte pas de mathématiques.}

$A$ \uncertain{ens.} : $S$, \ill{}, \ill{} les \uncertain{associés} constituent un
\uncertain{dodécaèdre combinatoire} (dont un \uncertain{icosaèdre}
\uncertain{combinatoire})

$F$ = ens.\ des faces = ens.\ des 12 structures polygonales sur $\mathfrak{S}$\note{lecture du symbole incertaine.}
les polygones combinatoires \uncertain{associés}.

\page{43}
\note{feuillet de calculs et de croquis, sans texte suivi ; on transcrit ce qui se lit.}

\marginal{\uncertain{petits clous pour \ill{} \ill{}}}

\drawing{en haut, un rectangle portant « 57 » et « 54 » ; un hexagone aux côtés fléchés ; un hexagone dont les sommets sont joints par des diagonales en traits pleins et en pointillés.}

\[
S = \tau \times \varepsilon, \qquad
A = \tau \times (\varepsilon \setminus \omega(\tau)), \qquad
\widetilde{S} = \tau \times \varepsilon \times \omega(\tau)
\]
\[
\tau \times (\varepsilon \setminus \omega(\tau)) \setminus \omega(\tau') \qquad
\tau\times\varepsilon \cdot \tau\times\varepsilon' \qquad
\tau \times \varepsilon \times \varepsilon'
\]

\drawing{un schéma de flèches : $\tau \leftarrow \tau\times\varepsilon$ ; « Tritria » $\xrightarrow{2}$ « Hex » $\xleftarrow{3}$ « (tri)$'$ ».}

\drawing{dans la moitié inférieure, un grand diagramme de flèches entre les ensembles « (tria$'$) », « Hex », « Tritria », « (tria) $\times$ hex », « (tria $\times$ 2) $\times$ hex $+$ \ill{} », « hex $+$ diag », « hex $+$ sommets », « Hexorien $\simeq$ (tria$'$/Hex) », « (trilles) $= \tau\times\varepsilon\times\omega(\tau)$ », $S$, $R \rightrightarrows \tau\times\varepsilon,\ \tau\times\varepsilon'$, les flèches portant les degrés $2$ ou $3$ ; en bas, un hexagone avec ses diagonales.}

\page{44}
\note{feuillet de croquis ; on relève ce qui se lit.}

\drawing{à gauche, deux hexagones de sommets $x_1,\dots,x_6$, avec les paires $\{1,2\}$, $\{2,5\}$, $\{1,6\}$, $\{1,5\}$, $\{3,4\}$, $\{2,4\}$, $\{2,6\}$, $\{1,5\}$, $\{1,6\}$, $\{2,4\}$ écrites à côté.}

\[
15 \cdot 6 \cdot 4 \cdot 2 \cdot 1 = 720 = 6!
\]
\[
\{1\,3\}\,\{2\,4\}\,\{3\,5\}\,\{4\,1\}\,\{5\,2\} \qquad
\{1\,5\}\,\{3\,2\}\,\{5\,4\}\,\{2\,1\}\,\{4\,3\}
\]

\drawing{à droite, un pentagone $x_1x_2x_3x_4x_5$, la suite $\{1\,2\}\,\{3\,4\}\,\{1\,5\}\,\{2\,3\}\,\{4\,5\}$, les pentagones étiquetés $1,2,3,4,5$ dans deux ordres différents, et « $\{i, i+2\}$, $\{j, j+2\}$ ».}

\[
10\,P = 720, \qquad P = 72, \qquad P_0 = 12 = \text{\uncertain{nb des struct.\ \ill{}}}, \qquad 4!
\]
« l'un à 5 él. »

\page{46}
graphe discret : sans arêtes. Parties discrètes d'un graphe (discrètes pour
la structure \uncertain{induite}) \add{dont l'ens.\ des sommets} : doublets, triades\add{,
quadruplets, pentuplets, hexaplets}, ; parties discrètes de card $2,3$.

\struck{graphe} graphe \struck{\ill{}} bidiscret : graphe $\Gamma$, décomposé en deux parties discrètes
$A$, $A'$ ($\Gamma = A \amalg A'$) \struck{tel que $\forall x \in A$ $\exists!\, x' \in A'$ tel que} \add{la relat.\ $R$ entre $A$, $A'$ définie par
$\{x, x'\} \in R \Leftrightarrow \{x,x'\}$ est un doublet (i.e.\ $x$, $x'$ non liés) est bijective}
\uncertain{sauf} un doublet. Cette décomposition est unique sauf si
$\operatorname{card} A = 2$ \add{i.e.\ $\operatorname{card}\Gamma = 4$} (les graphes de ce dernier type sont tous isom.\ à
$A_2 \amalg A_2$, i.e.\ somme disjointe de deux segments --- et le
choix d'une \struck{$\Delta$} décomposition \add{de $\Gamma$} en $A$, $A'$ équivaut à le
\uncertain{donnée} d'une bijection entre ces segments. Bi-doublets,
bi-quadruplets, bipentuplets, bihexaplets…
bitriades… \struck{parties bi-\ill{} discrètes d'ordre $n$ (bi}

\drawing{en marge, à hauteur de « $\operatorname{card} A = 2$ » : deux segments croisés de sommets $a, b, a', b'$.}

Dans un graphe cubique, les bases et bibases
sont les \struck{systèmes $\Gamma = A, A'$} hexaplets et bihexaplets.

\uncertain{Deux} parties discrètes d'un graphe sont dites \uncertain{« W-liées »}\note{le mot est récurrent dans tout le lot : une capitale en forme de \emph{W} (ou \emph{V}) suivie de « -liées » ou « -iées » ; la lecture « W-liées » n'est qu'une approximation du tracé. On le garde tel quel partout.} \struck{\ill{}}
si elles sont disjointes et si leur réunion est un \struck{\ill{}} \add{bidiscret par sa décomposition naturelle} graphe \uncertain{(\ill{} \ill{})}.
\textbf{NB} La catégorie des graphes bi-discrets d'ordre $n$ (card
\ill{}) équivalente canoniquement à la catégorie $(\mathrm{Ens}_2) \times (\mathrm{Ens})_n$.

Deux triades $\tau = \{a,b,c\}$, $\tau' = \{x,y,z\}$ sont \struck{dites} \add{dites} liées
si \struck{\ill{}} elles sont disjointes et leur réunion
est discrète, i.e.\ ssi leur réunion est un hexaplet \ill{}
ou encore une base.
\struck{Elles sont dites \uncertain{bi-liées}}

\textbf{Proposition} Soit $\Delta = (S, A)$ un graphe cubique, soit
$(\mathrm{Tria})$ l'ens.\ de ses triades, considérons sur $(\mathrm{Tri})$ la structure
de graphe définie par la relation des « triades liées »
\struck{\ill{}} \add{($\Delta$ ou au moins \ill{})}\note{insertion interlinéaire très effacée, lecture approximative.}
\uncertain{précédente}. Alors (1) $\mathrm{Tri}$ est \struck{\ill{} disjoint} \add{un graphe somme} de
\add{i.e.\ un \uncertain{cycle} cocaroé \uncertain{unit}…} \add{des hexagones} \struck{hexagones}. (2) Deux triades sont
\add{W-liées au sens plus haut} \struck{biliées}. ssi elles sont
\struck{W-liées} \add{W-liées} \struck{\ill{}} sur un même hexagone, et ssi elles sont
(dans ce cas $\tau, \tau'$) (3) Deux triades \add{$\tau, \tau'$} sont
\uncertain{opposées} dans un même hexagone \ill{}
\uncertain{la triade} $\exists!$ triade W-liée) (3) Deux triades sont
biliées ssi elles sont W-liées : \ill{} $=$ hexagone,

\marginal{\uncertain{card} $\mathrm{Tri} = 720$ …}
\marginal{NB les bitriades \uncertain{sont exactement} les hexagones}
\marginal{Mais sont distinctes liées ssi elles \uncertain{sont} des parties \ill{} et \ill{} de bases \ill{} \ill{} \ill{} \ill{} \ill{} \ill{} \ill{} \uncertain{distinctes}. Elles \struck{\ill{}} \ill{} \ill{} \ill{} il \ill{} \ill{} (dans \ill{} \uncertain{disjointes}) \ill{}}
\marginal{\struck{hex.\ gr.} \struck{hexagones} \struck{des triades} des triades : \uncertain{des quatre triades} \struck{W-liées} \add{hexagones} \struck{des triades} \ill{} ; $\operatorname{card} \mathrm{Hex} = 120$}
\note{la marge gauche, écrite de biais, est très serrée ; les quatre notes marginales ci-dessus en donnent ce qui se lit, sans ordre assuré. La numérotation des points de la proposition (deux fois (3)) est celle du feuillet ; la phrase sur les triades « opposées » est en partie couverte de ratures.}

\page{47}
(3)… est ni égale, ni liée, ni W-liée, i.e.\ ssi \add{$\tau \neq \tau'$ et} il existe
une triade $\tau''$ liée à $\tau$ et $\tau'$. (4) Soit $\tau$ une triade,
soient $\tau'$, $\tau''$ les deux triades liées à $\tau$. Alors
\begin{enumerate}
\item[a)] $\tau' \amalg \tau'' = L(\tau) = \{\text{ens.\ des sommets liés à aucun} \in \tau\}$
\end{enumerate}
\struck{\ill{}} En particulier (compte tenu de (3)) $L(\tau)$ est une
bitriade, et $L(\tau) = \tau' \cup \tau''$ est sa décomposition canonique
(ce qui explique la construction des deux triades $\tau', \tau''$
contiguës à une triade $\tau$) (5) Soit $\overline{\tau} = E(\tau)$ l'ens.\
des sommets liés à tous les él.\ de $\tau$. Alors $\overline{\tau}$ est une
triade (dans le triade « opposée » à $\tau$). D'où (6) les deux
triades biliées à $\tau$ s'obtiennent ainsi :
\[
\begin{aligned}
&\text{l'une est} && \overline{\tau'} = L(\tau'') \setminus \tau\\
&\text{l'autre} && \overline{\tau''} = L(\tau') \setminus \tau
\end{aligned}
\]
\struck{(7) des hexagones des triades} \add{L'hexagone}
\uncertain{des} triades contenant $\tau$ est \struck{\ill{}}
$\tau\ \tau'\ \overline{\tau''}\ \overline{\tau}\ \overline{\tau'}$ (la suite est des côtés
circulaires…) \struck{\ill{}}. Aussi \ill{} \ill{}

\begin{tikzcd}[column sep=small, row sep=small]
 & \tau \arrow[rr, no head] & & \tau'' \arrow[dr, no head] & \\
\tau' \arrow[ur, no head] & & & & \overline{\tau'} \arrow[dl, no head] \\
 & \overline{\tau''} \arrow[ul, no head] \arrow[rr, no head] & & \overline{\tau} &
\end{tikzcd}

\note{l'hexagone du feuillet ; deux de ses côtés sont doublés ou fléchés, ce que le dessin ci-dessus ne rend pas.}

(7) Pour la relation de « biliées », \struck{les triades} \add{liées} $(\mathrm{Tri})$
\uncertain{l'ens.}\ des triades est un graphe somme de
triangles (dits \struck{triangle} \struck{triples} « triangles » de triades biliées), dans
l'ens.\ des triangles de triades biliées \struck{\ill{}} le groupe
\ill{} « W-liés » est une involution \uncertain{à} pt fixe.
Les hexagones correspondent aux orbites de cette involution,
\uncertain{les triangles} correspondant \struck{\ill{}} à \struck{une} triangle et \ill{} (de cardinal $6.3 = 18$) \add{\uncertain{liées}} des triangles
(8) La réunion \ill{}
d'un hexagone \ill{} a un complémentaire
(de cardinal $27 - 18 = 9$) qui est une bitriangle.
L'application qui à tout hexagone de
triades associe \struck{\ill{}} bitriangle précédent,
est bijective. \struck{Les hexagones \ill{}}

\marginal{Tritria = triangles de triades biliées ; cardinal $2.120 = 240$}

\drawing{dans la marge inférieure gauche, un hexagone de sommets $\tau_1$, $\tau_2$, $\tau_3$, $\overline{\tau_1}$, $\overline{\tau_2}$, $\overline{\tau_3}$, avec ses diagonales en pointillés ; le feuillet est barré de trois traits obliques sur sa partie basse.}
\note{la fin du point (8) est barrée de longs traits obliques ; on la donne néanmoins.}

\page{48}
(8) Soit \struck{\ill{}} $\{\tau_i\}_{i\in I}$ \struck{\ill{}} un triangle de triades liées \struck{Alors}
\add{(\uncertain{relation} \uncertain{non-liées} \uncertain{Notat.} \ill{})}
soit $I = \{i,j,k\}$,
\struck{Soit} Si $x_i \in \tau_i$, et si $x_j$ ($x_k$) est l'unique él.\ de
$\tau_j$ ($\tau_k$) non lié à $x_i$, alors $x_j$ et $x_k$ sont non liés.
En d'autres termes, si $\Gamma \subset \prod_{i\in I}\tau_i$ \struck{\ill{}} est l'ens.\
des $(x_i)_{i\in I}$ \struck{tels} tels que $x_i$ non lié à $x_j$ pour $i \neq j$),
alors les projections $\Gamma \xrightarrow{p_i} \tau_i$ sont bijectives (donc
\ill{} \ill{} un système transitif d'isom.\ entre les $\tau_i$, $i\in I$),

\struck{(9) Soit \ill{} $X = \tau_1 \cup \tau_2 \cup \tau_3$ la réunion des trois triades d'un
triangle $\{\tau_1, \tau_2, \tau_3\}$ des triades biliées [$\overline{X} = \overline{\tau_1} \cup \overline{\tau_2} \cup \overline{\tau_3}$
la réunion des trois triades du triangle W-lié,
\ill{} $X$, $\overline{X}$ sont de cardinal 9 et disjoints].}

(9) Soit $X = \coprod_{i\in I} \tau_i \hookrightarrow S$, \add{(card $X = 9$)} d'où $X \to I$, dont
les fibres sont \struck{\ill{}} \add{les} 3 triades du triangle \uncertain{des triades donné},
et soit $X \to J$ l'application dont les fibres \uncertain{sont les}
\add{$(\tau_j')_{j\in J}$} \add{donnant les fibres} \uncertain{sont} \uncertain{des} $3$
\uncertain{autres} \struck{triades} \struck{\ill{}} triades \struck{\ill{}} disjointes. Ceci posé

($\cdot$) $X$ est une bitriangle, \struck{dont} \add{dont} les six \add{triades} \struck{\ill{}}
(ou \uncertain{antitriangles}) sont les $\tau_i$ \struck{\ill{}} ($i \in I$) $\tau_j'$ ($j\in J$)) formant
les 2 familles \struck{de triades} \add{\uncertain{\ill{}}} de $X$. Aussi
\struck{l'application $X \to I \times J$ est bijective}
$I$, $J$ s'identifient resp.\ à l'une et l'autre
famille de triades du bitriangle $X$, et l'application
\[
X \longrightarrow I \times J
\]
est bijective.

(10) La donnée d'un triangle de triades biliées \add{$(\tau_i)_{i\in I}$}
équivaut à celle d'un bitriangle $X$, muni
\struck{\ill{}} du choix d'une \struck{\ill{}} des deux familles \add{$I$, $J$} de
\add{triades} \uncertain{mutuellement} disjointes de $X$ (\uncertain{coorientation})

\page{49}
\struck{\textbf{NB} \ill{} pour que deux triades $\tau_1, \tau_2$ \ill{} \ill{} biliées,
il faut et il suffit qu'elles soient disjointes et appartiennent à un
bitriangle, qui est unique et déterminé, et est aussi
l'unique \ill{}}\note{passage biffé de traits obliques.}

(11) $X$ est l'unique bitriangle contenant $\tau_i$ --- deux triades
\uncertain{non} biliées ssi elles \uncertain{sont} \add{disjointes et,} contenues dans, un \ill{}
bitriangle \struck{\ill{}}

(12) Soit \struck{\ill{}} $(\tau_i)_{i\in I}$ un triangle de triades biliées,
\uncertain{réunion} du bitriangle $X = \bigcup_{i\in I} \tau_i \subset S$ \ill{} \ill{}
coorientation $\{\tau_i\}$ \struck{\ill{}}, soit \add{$(\tau_j')_{j\in J}$} \struck{$(\tau_i)_{i\in I}$ \ill{} le triangle des}
\struck{triades} de $X$ de l'autre famille (de sorte que $X \xrightarrow{\sim} I \times J$).
Soit $(\overline{\tau_i})_{i\in I}$ le triangle de triades « W-lié », $\overline{X} = \bigcup_{i\in I}\overline{\tau_i}$,
et $\overline{J}$ l'ens.\ des triades de $\overline{X}$ de l'autre famille,
donc $\overline{X} \xrightarrow{\sim} I \times \overline{J}$. Soit $\widetilde{X} = S \setminus X \setminus \overline{X}$, qui
est de cardinal $9$. \add{$27 - 9 - 9$} Alors
\begin{enumerate}
\item[a)] $\widetilde{X}$ est un \struck{bitriangle} bitriangle. \struck{Soit} L'application qui à tout
hexagone de triades \uncertain{associe} $\widetilde{X}$, est bijective. \struck{\ill{}}
\end{enumerate}

\page{50}
\struck{b) L'ens.\ des coorientations de} $\widetilde{G}$ d'automorphismes de $X$,
b) \struck{Soit $\widetilde{X}$ \ill{} les deux \ill{} bitriangles
disjoints de $X$ \ill{} \ill{} \ill{}} Considérons le groupe \struck{\ill{}} \add{\uncertain{formé}} par les
\add{\ill{} \ill{}} des symétries
\struck{\ill{}} $\sigma_t$ ($t \in \mathrm{Tri}(\widetilde{X})$). Alors
  \begin{enumerate}
  \item[$\alpha$)] $\widetilde{G} \longrightarrow \operatorname{Aut}_{\mathrm{gr}}(\widetilde{X})$ \add{\uncertain{induit}} est un isom.\ \add{de $\widetilde{G}$} sur le groupe des
  automorphismes \uncertain{intérieurs} de $\widetilde{X}$
  \item[$\beta$)] Les orbites de $\widetilde{G}$ opérant sur $S \setminus \widetilde{X}$ $(= X \amalg \overline{X})$
  sont de cardinal $6$ --- donc il y a \uncertain{trois} telles
  orbites [\uncertain{qui sont les réunions de la \ill{} forme} $\tau \cup \overline{\tau}$ (triade de \uncertain{bihexaH})]
  [\uncertain{On a aussi} $\widetilde{G} \simeq \mathfrak{S}_{\mathbb{F}} \times \mathfrak{S}_{\mathbb{Q}}$, où $\mathbb{F}$, $\mathbb{Q}$ sont
  les deux ens.\ de card.\ $3$ des triangles de $\widetilde{X}$ de
  l'une et l'autre espèce, \struck{les \uncertain{orbites} sont} $\mathfrak{S}_{\mathbb{F}}$ et $\mathfrak{S}_{\mathbb{Q}}$
  \ill{} simplement transitifs sur chacune des
  trois orbites] \add{Dans l'ens.\ des trois orbites s'identifie : …
  \ill{} $I$ … des 3 diagonales de $H$.}
  \end{enumerate}

c) \struck{\ill{}} \add{Si $t \in \mathrm{Tri}(\widetilde{X})$, $\sigma_t$ échange $\tau$ et $\overline{\tau}$}
\add{$\widetilde{G}^{\circ} \subset \widetilde{G}$ formé par} les \uncertain{produits pairs} de
dans le groupe \struck{\ill{}} \ill{}
symétries \struck{\ill{}} \uncertain{associées} à chacun des $\tau \in H$ ---
et les orbites de $\widetilde{G}^{\circ}$ sur $S \setminus \widetilde{X}$ sont de cardinal
$3$, et ce sont les $6$ triades de l'hexagone $H$.
\struck{(\ill{} l'ens.\ des \ill{} \ill{}}
\struck{d) Soit}

(11) Soient $(X_\alpha)_{\alpha\in A}$ une famille de trois bitriangles
mutuellement disjoints. Pour $\alpha \in A$, soit $I_\alpha$ l'ens.\
de cardinal $3$ associé à $X_\alpha$ \add{l'ens.} \struck{\ill{}} \uncertain{trois} orbites
du $\operatorname{Aut}_{I_\alpha}(X_\alpha)$ opérant sur $S \setminus X_\alpha = \bigcup_{\beta\in A\setminus\{\alpha\}} X_\beta$,
\uncertain{donc} \ill{} \ill{} une application $X_\beta \xrightarrow{p_{\alpha\beta}} I_\alpha$ ($\beta \neq \alpha$)
\note{la numérotation reprend (11), déjà employé p.~49 ; elle est celle du feuillet.}

\page{51}
\uncertain{de} degré 3, dont les \add{3} fibres sont des \add{trois} triades
(\uncertain{mutuellement} disjointes) formant l'une des deux
familles de triades de $X_\beta$. Alors
\begin{enumerate}
\item[a)] \struck{Les \ill{}} \struck{deux} \uncertain{Les deux} familles de triades, fibres des
deux applications $X_\alpha \to I_\beta$, $X_\alpha \to I_\gamma$ (où \struck{\ill{}}
$A = \{\alpha,\beta,\gamma\}$) sont distinctes, donc l'application
correspondante $X_\alpha \xrightarrow{\sim} I_\beta \times I_\gamma$ est bijective.
\end{enumerate}
Plus intrinsèquement, $\forall \alpha \in A$, on a
\[
X_\alpha \xrightarrow{\sim} \prod_{\beta \in A\setminus\{\alpha\}} I_\beta
\]
d'où une bijection naturelle
\[
S \simeq \coprod_{\alpha} \prod_{\beta\in A\setminus\{\alpha\}} I_\beta \quad \text{\struck{$I_\alpha$}}
\]
dans le cas où $A = \{1,2,3\}$, cela fait
\[
S \simeq \underbrace{I_2 \times I_3}_{X_1} \amalg \underbrace{I_3 \times I_1}_{X_2} \amalg \underbrace{I_1 \times I_2}_{X_3}
\]
On a $\operatorname{coor}(X_\alpha) \simeq{}$ \ill{} $A \setminus \{\alpha\}$
\begin{enumerate}
\item[c)] Pour que \struck{$x_\alpha \in X_\alpha$, $y_\alpha \in$} \add{\ill{}} $x_\alpha, y_\alpha \in X_\alpha$ \struck{\ill{}},
soient non liés il f.\ et s.\ que \struck{$\exists$} $\beta \in A \setminus\{\alpha\}$ \struck{\ill{}}
tel que $p_{\beta\alpha}(x_\alpha) = p_\beta(y_\alpha)$ (si $x_\alpha \neq y_\alpha$, ce $\beta$ est
unique) --- pour qu'ils soient liés, il f.\ et s.
que $\forall \beta \in A \setminus \{\alpha\}$, on ait $p_{\beta\alpha}(x_\alpha) \neq p_{\beta\alpha}(y_\alpha)$
\item[d)] Pour que $x_\alpha \in X_\alpha$ et $y_\beta \in X_\beta$ ($\alpha\neq\beta$) soient liés
il f.\ et s.\ que (si $\gamma$ \struck{\ill{}} \struck{\ill{}} est l'\struck{\ill{}} \add{él.} $\neq \alpha,\beta$ de $A$) on ait
\add{$p_{\gamma\alpha}(x_\alpha) = p_{\gamma\beta}(x_\beta)$}
\end{enumerate}

\page{52}
$S \simeq$ ens.\ des parties \add{$A$} de cardinal $2$ de $I$,
telles que $p|A$ soit injectif (contenu parallèles, \uncertain{ou}
partition de cardinal 2)
\[
X_\alpha = \text{ens.\ des } A \text{ tels que } p(A) = \Lambda \setminus \{\alpha\}
\]
\[
A, B \in S \text{ liés ssi } A \neq B \text{ et}
\]
\[
\begin{aligned}
&\text{ou bien } p(A) = p(B), && A \cap B = \emptyset\\
&\text{ou bien } p(A) \neq p(B), && A \cap B \neq \emptyset
\end{aligned}
\]

\begin{tikzcd}[row sep=small]
I \arrow[d, "p"] & I_\alpha \\
\Lambda & \alpha
\end{tikzcd}

\drawing{en marge gauche, trois rangées de trois points, l'une barrée d'un trait, avec une flèche verticale : le schéma $I \to \Lambda$.}

(12) La catégorie des graphes cubiques, munis d'une
décomposition en trois \struck{bitriangles} \add{bitriangles} disjoints, est
équivalente : \struck{\ill{}} \add{à la cat.} \struck{(des ens.\ de type} $(\mathrm{Ens})_{(3,3,3)}$, munis
d'une partition de type $(3,3,3)$. Le groupe
des automorphismes d'une telle structure est
\[
\simeq \mathfrak{S}_3 \cdot (\mathfrak{S}_3 \times \mathfrak{S}_3 \times \mathfrak{S}_3)
\]
qui est d'ordre $\uncertain{3 \cdot 6^4}$ (c'est donc \add{presque} un \uncertain{sous-groupe} de Sylow
\add{\uncertain{il} \ill{} \ill{} des \struck{\ill{}} \add{cinq} dans les $I_\alpha$ et
dans $\Lambda$ et exiger qu'elles soient compatibles…}
\ill{} de $\operatorname{Aut}(\Delta)$ \struck{$\ill{}$} !) La cat.\ des couples
$(\Delta, X)$ d'un graphe cubique, et d'un
bitriangle $X$ dans $\Delta$, est équivalente : \add{la cat.} \struck{\ill{}},
\struck{complexe} \uncertain{de} $(\mathrm{Ens})_3 \times (\mathrm{Ens})_{3,3}$ des couples formés
d'un ens.\ de card.\ $3$ (savoir l'ens.\ \struck{\ill{}} \add{\uncertain{des \ill{}}}
\uncertain{orbites} associé : $X$ \uncertain{en} $\Delta$) et d'un ens.\ \uncertain{à partition}
de type $(3,3)$ (savoir l'un des triades de $X$,)
\uncertain{avec} sa partition naturelle).

\marginal{NB L'\uncertain{ordre} \struck{\ill{}} $= 3^4 2^4$
de \uncertain{Sym} \uncertain{la} $X_\alpha$
se réunissent \uncertain{dans}
de ce sous-groupe
de $\operatorname{Aut}(\Delta)$,
commun l'un
des $3$ orbites sur ce
groupe : dans ce \uncertain{bi}-
\uncertain{sous-groupe} de Sylow
est \uncertain{dans} \uncertain{presque}
\uncertain{normalisateur}.}
\note{la note marginale, écrite de biais en lignes très courtes, n'est lue que par fragments.}

\page{53}
(13) L'ens.\ des 18 triades \struck{de} $S$ contenues dans un
des $X_\alpha$ est \struck{en} conv.\ binunivoque avec \struck{l'ens.} le
sous-ens.\ $T$ de $I \times \Lambda$ formé des couples
$(i, \alpha)$ avec $p(i) \neq \alpha$, en associant : à un tel
couple la triade \add{$\tau_{i,\alpha} \subset X_\alpha$} (formée des $A \in S \subset \mathfrak{P}_2(I)$) tels
que $p(A) = \Lambda \setminus \{\alpha\}$, $A \ni i$. \add{\uncertain{distincts}} La triade W-liée
à $\tau_{i\alpha}$ est $\tau_{i\beta}$, où \struck{$\beta$} \uncertain{$\beta$} \struck{\ill{}} \add{l'autre} él.\ de \struck{$\Lambda$}
$\Lambda \setminus \{p(i)\}$. Les deux triades \struck{\ill{}} liées à $\tau_{i\alpha}$ sont
$\tau_{j,\beta}$, $\tau_{k,\beta}$, où $j$, $k$ sont les \add{2 autres} \struck{\ill{}} él.\ de \struck{$p^{-1}(\ill{})$} $I$
au-dessus de $p(i)$. Les deux triades biliées à $\tau_{i\alpha}$
sont $\tau_{j,\beta}$\struck{,$\ill{}$}, $\tau_{k,\beta}$. L'hexagone formé des
six triades est formé de
l'ens.\ des $\tau_{i,\alpha}$ tels que $p(i) \neq \alpha$
soit donné. L'ens.\ de ces hex.\
est donc en corr.\ 1-1 avec $\Lambda$
\struck{\ill{} \ill{} \ill{}} [à l'hexagone $H$ correspond
le complémentaire de sa réunion dans $S$]

\drawing{à gauche, un graphe : les sommets $p, q, r$ (en colonne) et $i, j, k$ (en colonne), joints par des traits pleins ($p$–$i$, $q$–$i$, $r$–$i$) et pointillés, avec une troisième colonne de points à droite reliés à $i$ par des pointillés.}

(14) \struck{L'}L'ens.\ \add{à 6 él.}\ des couples formés d'un des \struck{\ill{}} bitriangles
$X_\alpha$, et d'une coorientation de $X_\alpha$, est $\simeq$
l'ens.\ des \struck{\ill{}} couples $(\alpha, \beta) \in \Lambda^2$ avec
$\beta \neq \alpha$ (i.e.\ : l'ens.\ des repères $(\alpha,\beta,\gamma)$ de
l'ens.\ $\Lambda$ de cardinal 3), à $(\alpha,\beta)$ correspond…
\note{les quatre dernières lignes sont barrées de trois traits obliques.}

\page{54}
le couple $(X_\alpha, \{\tau_i$
\note{seule ligne du feuillet, inachevée.}

\page{55}
Sa structure combinatoire est donnée par les
ens.\ $I_\gamma$ de cardinal 3 et \struck{\ill{}} $\Lambda_\gamma = \Lambda \setminus \{\gamma\}$ \add{$= \{\alpha,\beta\}$} de cardinal 2,
comme
\[
H \simeq I_\gamma \times \Lambda_\gamma,
\]
où $\tau_{i,\alpha}$ et $\tau_{j\beta}$ \uncertain{liés} ssi $i \neq j$, $\alpha\neq\beta$. \struck{\ill{} \ill{}}

(14) Les triangles de triades biliées
correspondent aux couples
$(\alpha, \beta)$ d'él.\ distincts de $\Lambda$
(i.e.\ aux repères de $\Lambda$) en
associant à tel repère \struck{\ill{} \ill{}} \uncertain{tel} \ill{}
triades $T_{\alpha,\beta,\gamma}$ \struck{$\{\tau_{i,\alpha}, \tau_j\}$} $\{\tau_{i\alpha}\}_{i \in I_\gamma}$. Cet ens.\ \ill{}
\uncertain{W-lié} \ill{} \ill{}

\begin{tikzcd}[column sep=small, row sep=small]
 & \tau_{i\alpha} \arrow[r, no head] & \tau_{j\beta} \arrow[dr, no head] & \\
\tau_{k\beta} \arrow[ur, no head] & & & \tau_{k\alpha} \arrow[dl, no head] \\
 & \tau_{j\alpha} \arrow[ul, no head] \arrow[r, no head] & \tau_{i\beta} &
\end{tikzcd}

\[
\overline{T_{\alpha\beta\gamma}} = T_{\beta,\alpha,\gamma}
\]
Le triangle de triades $\widetilde{T_{\alpha\beta\gamma}}$ \struck{\ill{}} contenu
dans le bitriangle $X_\alpha$ \uncertain{que} \struck{\ill{}} $T_{\alpha\beta\gamma}$,
et distinct de $T_{\alpha\beta}$, est $T_{\alpha\gamma\beta}$. Donc
le groupe engendré par les \uncertain{opérations} --- et
$\sim$ sur l'ens.\ des triangles de triades liées
\note{la ligne médiane de la page, entre le diagramme et la formule, n'est pas lue.}

\page{56}
d'un graphe cubique $\Delta$ est isomorphe au
groupe $\mathfrak{S}_3$, où $-$ correspond à $\sigma_{12}$,
$\sim$ à $\sigma_{23}$ …

\textbf{NB} \uncertain{La donnée} d'un couple $(X, \{\tau_i\}_{i\in I})$
d'un bitriangle $X \subset \Delta$ et d'une coorientation
de celui-ci, \uncertain{donc} \uncertain{presque} \uncertain{associée}
à un tel couple, \uncertain{d'un} \uncertain{autre} \uncertain{tel}
\struck{couple $(X', (\tau_i')_{i\in I})$ formé du bitriangle $X'$ \add{la réunion} \uncertain{associé} à $X$, tel que les $\tau_i$ soient}
\struck{les orbites} $X' = \bigcup_{i\in I} \overline{\tau_i}$ \add{\uncertain{réunion} de}
\add{$\hookrightarrow$ triades W-liées}
$(\overline{\tau_i})_{i\in I}$, d'autre part. $X$ muni de
l'autre coorientation $(\tau_i')_{i\in I'}$ …

\struck{(15)} \struck{\ill{}}
(15) Les triangles de $X_\alpha$ correspondant aux
\uncertain{fonctions} $t \subset \prod_{\beta\in\Lambda\setminus\{\alpha\}} I_\beta$ \struck{qui sont des} telles que
$t \to I_\beta$ ($\beta \in \Lambda\setminus\{\alpha\}$) soit
\add{i.e.\ (si $\Lambda = \{\alpha,\beta,\gamma\}$) aux bijections $I_\beta \xrightarrow{\sim} I_\gamma$}
les applications… \uncertain{telles} \struck{\ill{}}
bijections, en associant : une
\struck{l'ens.\ des \ill{}} \add{bijection} $\varphi$ : \struck{\ill{}} l'ens.\ des él.\ $\{x_\beta, \varphi(x_\beta)\}_{x_\beta\in I_\beta}$
de $X_\alpha$. Ceci posé, pour un tel triangle
$t = t_\varphi$, la symétrie correspondante de $\Delta$ est

\page{57}
telle associée : l'automorphisme de
$(I \xrightarrow{p} \Lambda)$ : son action sur $\Lambda$ est la
permutation qui fixe $\alpha$ et échange $\beta, \gamma$,
\ill{}\note{le milieu du feuillet n'a pas été lu de près ; on en donne le début et la fin.}
\begin{enumerate}
\item[b)] opérant trivialement sur l'ens.\ $\pi(X_\alpha) \simeq I_\alpha$
des « trois orbites » associé à $X_\alpha$.
\end{enumerate}

(16) Soit $X \subset S$ un bitriangle dans $S$, et
\uncertain{\ill{}} l'un des « trois orbites » de $S \setminus X$.
\add{$\{\mathcal{X}_i\}_{i\in I_X}$} \add{Si} L'un des « trois orbites » de $S\setminus X$.
Alors les \uncertain{sous-graphes} \add{de type} $\mathcal{T}_1$ \struck{\ill{}} contenant
$X$ sont les $X \cup \mathcal{X}_i$. Donc la \add{catégorie} \struck{\ill{}}
\uncertain{des} triplets $(S \supset \mathcal{X} \supset X)$ d'un \add{graphe} \struck{\ill{}}

\page{58}
cubique \add{($\mathcal{T}_2$)} muni d'un $\mathcal{T}_1$ et d'un $\mathcal{T}_0$ contenu
dans celui-ci, équivaut : la \add{catégorie} \struck{\ill{}}
\struck{des} $(\mathrm{Ens})_{3,3} \times (\mathrm{Ens})_2$ \struck{[en associant l'ens.\ des \ill{}} \struck{\ill{} associés à un}
\struck{ens.\ à trois éléments],}
triades de $X$, et l'ens.\ des deux orbites dans $S \setminus \mathcal{X}$]
\struck{Complexe de} La \add{catégorie} \struck{\ill{}} des systèmes
$(S \supset \mathcal{X} \supset X \supset t)$ d'un $\mathcal{T}_2$ contenant un $\mathcal{T}_1$ contenant un
$\mathcal{T}_0$ contenant un $\mathcal{T}_{-1}$, équivaut : la catégorie
$(\mathrm{Ens})_3 \times (\mathrm{Ens})_2 \times (\mathrm{Ens})_2$, [en lui associant l'ens.\
des points de l'une des triades de \struck{$\mathcal{X}$} $X$ \uncertain{par} la
relation d'équiv.\ définie par $t$, l'ens.\ des
deux coorientations de $X$, et l'ens.\ des deux
orbites de $S \setminus \mathcal{X}$] …

(17) Soit $\mathcal{X}$ un $\mathcal{T}_1$ dans \uncertain{le} $\mathcal{T}_2$ $S$. Alors
$S \setminus \mathcal{X}$ est une bibase de $S$, et on trouve
ainsi une bijection entre l'ens.\ des
$\mathcal{T}_1$ de $S$, et l'ens.\ des bibases. La
catégorie des $\mathcal{T}_2$ munis d'une bibase, celle des
$\mathcal{T}_2$ munis d'un $\mathcal{X}_1$, et la catégorie des
\ill{} \ill{} sont \uncertain{équivalentes}.

\page{59}
(18) À tout $B \in (\mathrm{Ens})_6$, ensemble de cardinal $6$,
associons le graphe $\mathcal{T}_1(B)$ suivant
\[
\begin{cases}
\text{pour sommets } = \mathfrak{P}_2(B)\\
\text{\struck{$A$, $B$}} \ u, v \in \mathfrak{P}_2(B) \text{ liés ssi } u \cap v = \emptyset
\end{cases}
\]
Alors $\mathcal{T}_1(B)$ est un graphe \uncertain{tritétraèdre} \uncertain{donc}
un $\mathcal{T}_1$, et le foncteur $B \mapsto \mathcal{T}_1(B)$ est
une équivalence de la catégorie $(\mathrm{Ens})_6$ avec
la catégorie des $\mathcal{T}_1$. Le foncteur quasi-inverse
\uncertain{associe} : $\mathcal{X}$ l'ens.\ $B$ des
parties discrètes de card $5$ de $\mathcal{X}$. Alors \ill{}
$x \in \mathcal{X}$ est contenu dans exactement deux
éléments de $B$, et l'application $\mathcal{X} \to \mathfrak{P}_2(B)$
ainsi obtenue est un isom.\ du graphe
$\mathcal{X}$ avec le graphe $\mathcal{T}_1(B)$. La donnée d'un
bitriangle dans $\mathcal{X}$ équivaut : la donnée
d'une partition de type $3,3$ de $B$
--- à celle-ci correspond l'ens.\ des \struck{\ill{}} \add{($\subset \mathfrak{P}_2(B)$)}
qui sont des sections de $B$ sur l'ens.\ quotient
à deux éléments. Tout bitriangle a un
$\mathcal{T}_2$ \uncertain{multiple} canonique (le foncteur
des $(\mathcal{X}, X)$ d'un $\mathcal{T}_1$ et d'un $\mathcal{T}_0$ contenu,
vers les $X$, est une équivalence de catégories).

\page{60}
19) Considérons \struck{\ill{}} \add{graphe} cubique $\Delta$ avec une partition
\uncertain{tritétraèdre} $S = \coprod_{\alpha\in\Lambda} X_\alpha$, et considérons le groupe
$G\ (\simeq \mathfrak{S}_3 \cdot (\mathfrak{S}_3 \times \mathfrak{S}_3 \times \mathfrak{S}_3))$ des automorphismes de
\uncertain{Ceci} \uncertain{dit}, \struck{\ill{}} et \uncertain{une} \add{\uncertain{équivalence} des} \add{\uncertain{il y a} l'ens.\ des (bitri) des bitriangles de $\Delta$.}
\add{$(\mathrm{Tria})$} l'ens.\ des triades de $\Delta$. Alors \struck{il y a}
a) \struck{Il y a} \add{ex.\ \uncertain{tout}} \add{\uncertain{quatre}} \struck{trois} orbites de $G$ sur $\mathrm{Tria}$ :
  \begin{enumerate}
  \item[1°)] L'orbite formée des triades contenues
  \ill{} dans un des $X_\alpha$
  \item[2°)] L'orbite \struck{formée} des triades de la forme
  \drawing{trois colonnes de trois points ; deux traits partent d'un point de la colonne médiane vers les deux autres colonnes.} (type I)
  \item[3°)] L'orbite \struck{formée} des triades de la forme
  \drawing{deux traits issus d'un même point vers deux points de la colonne voisine.} (type II)
  \item[4°)] L'orbite des triades de la forme
  \drawing{trois traits horizontaux sur trois rangées de points.} (type III)
  \end{enumerate}

b) Les triades contiguës : une triade
de type II est de type I, celles
contiguës à une triade de type I sont
de type I l'une, de type II l'autre \add{(d'une triade de type I ou II)}
de sorte que \struck{\ill{}} l'hexagone \struck{\ill{}} \struck{triades} est
de type

\begin{tikzcd}[column sep=small, row sep=small]
 & \mathrm{I} \arrow[r, no head] & \mathrm{II} \arrow[dr, no head] & \\
\mathrm{II} \arrow[ur, no head] & & & \mathrm{II} \arrow[dl, no head] \\
 & \mathrm{II} \arrow[ul, no head] \arrow[r, no head] & \mathrm{I} &
\end{tikzcd}

(2 sommets \add{opposés} de
type I, les autres de type II), l'hexagone
\add{formé de triades} d'une triade de type (III), est \struck{de type III}…
\note{la phrase se poursuit au-delà du lot.}

\end{document}
